% ============================================================
%  THE ORTYX PROTOCOL
%  Part III — Decomposition
%  Chapter 12: Failure Geometry Analysis
% ============================================================

\chapter{Failure Geometry Analysis}
\label{ch:failure-geometry}

\chapepigraph{To name a failure mode is not to condemn a framework.
It is to specify where the framework's development must go
if it is to become what it claims to be.}{Flyxion}

\noindent
Chapter~11 introduced the audit operators and applied them
at the level of the Phoenix corpus as a whole. Chapter~12
performs the case-by-case analysis: examining each of the
four failure geometries in detail, identifying the specific
loci in the corpus where each applies, and distinguishing
the degrees of severity that matter for the decomposition
in the Interlude.

The analysis proceeds in order of increasing severity:
$\FailGeo{1}$ (Decorative Formalism) is the most benign,
because operationalization gaps can in principle be closed
by additional specification work. $\FailGeo{4}$ (Operational
Severance) is more serious, because it requires not just
specification but the development of measurement
infrastructure. $\FailGeo{2}$ (Definitional Closure)
is more serious still, because closing it requires
restructuring the framework's foundational choices.
$\FailGeo{3}$ (Semantic Universality Collapse) is the
most severe, because at its limit it prevents any
observation from bearing on the framework's claims.

\section{Failure Geometry 1: Decorative Formalism}
\label{sec:fg1-analysis}

\subsection{The Harmonic Alignment Score}

The clearest instance of $\FailGeo{1}$ in the Phoenix
corpus is the harmonic alignment score $H_i$ in the
Marine salience functional. As noted in Chapter~7,
this component is described as measuring whether the
current peak frequency is consistent with a harmonic
series but is not defined precisely enough to generate
a unique implementation.

The severity of this instance is low. $H_i$ is operationally
grounded in concept: one can compute a harmonic alignment
score for a given peak given a reference fundamental,
a definition of harmonic series, and a threshold for
what counts as alignment. The specification needed is:
(1)~how the reference fundamental is determined, (2)~what
the definition of harmonic series is (integer multiples,
near-integer multiples within some tolerance), and (3)~what
the alignment threshold is and whether it is fixed or
adaptive. These are engineering decisions that do not
require new theory; they require commitment.

\begin{auditprop}
\label{ap:fg1-hcf}
The harmonic alignment score $H_i$ in the Marine functional
exhibits $\FailGeo{1}$ at a low severity level. The
operational gap is closable by engineering specification
without requiring theoretical revision. The functional
form of the salience metric is sound; the under-specification
affects the parameterization, not the architecture.
\end{auditprop}

\subsection{The Emotional Amplitude Modulation}

A more significant instance of $\FailGeo{1}$ appears in
the \MEMeight{} emotional amplitude modulation:
\[
  A_i(x,t) = A_i^0 \exp(-\lambda_i t)(1 + \eta\, e(t)).
\]
The mathematical form is precise. The emotional state
function $e(t)$ is not. What is $e(t)$? In a system
that does not have direct access to physiological
measures of emotional state, $e(t)$ would need to be
inferred from signal properties. Which signal properties?
Through what procedure? What scale is $e(t)$ on, and
what value does it take in neutral or ambiguous contexts?

Unlike $H_i$, this operational gap is not easily closed
by engineering commitment alone. Operationalizing $e(t)$
requires a theory of what signals carry emotional content
and how emotional content is measured from those signals
--- a non-trivial research programme in its own right.
The decoration here is at a higher level: the formal
expression is well-formed, but its realizability
presupposes a substantial prior solution to the problem
of affective computing.

\section{Failure Geometry 2: Definitional Closure}
\label{sec:fg2-analysis}

\subsection{The Authenticity Tautology}

The authenticity concept's most significant instance of
$\FailGeo{2}$ can be stated precisely. The corpus makes
the following moves:

\begin{enumerate}
  \item Authentic audio is audio that faithfully preserves
  the temporal structure of the original acoustic event.
  \item Lossy compression disrupts temporal structure
  (measured by Marine jitter metrics).
  \item Therefore, lossy compression reduces the
  authenticity of audio.
\end{enumerate}

Step~(3) follows from steps~(1) and~(2) in combination
with the definitional move in step~(1). But step~(1)
is not an empirical observation about what listeners
regard as authentic; it is a definitional choice about
what authenticity means. The argument would be equally
valid if step~(1) read: ``Authentic audio is audio
that faithfully preserves the spectral envelope of
the original'' --- in which case the conclusion about
lossy compression would not follow, because modern
codecs preserve spectral envelopes reasonably well.

The choice to define authenticity in terms of temporal
structure rather than spectral envelope determines the
conclusion. The conclusion cannot therefore serve as
evidence for the choice. But the corpus uses the
apparent fit between the definition and the engineering
results (the Marine detects what we care about: it
detects \emph{authenticity}) as though it were
confirmation of the framework rather than a consequence
of the definition.

\begin{auditprop}
\label{ap:fg2-authenticity}
The Phoenix corpus's authenticity claims exhibit
$\FailGeo{2}$ in the following specific form: the
concept of authenticity is implicitly defined in terms
of the Marine metric's primary variable (temporal
coherence), and the claim that the Marine detects
authentic audio is therefore not a discovery but a
definitional consequence. The severity is moderate:
the framework could escape $\FailGeo{2}$ by providing
an independent characterization of authenticity ---
one that predates and motivates the Marine metric
rather than following from it --- against which the
Marine's performance could be evaluated.
\end{auditprop}

\subsection{Meaningful Structure as Marine Output}

A related instance concerns the concept of ``meaningful
structure.'' The corpus repeatedly treats high-salience
Marine output as equivalent to meaningful signal content.
If meaningful structure is defined as what the Marine
detects, then the claim that the Marine detects meaningful
structure is trivially true. But the interesting claim
--- the one that would provide genuine support for
the Phoenix paradigm --- is that what the Marine detects
corresponds to some independently characterized notion
of meaningfulness. Establishing this requires either
(i)~a prior, independent operationalization of meaningful
structure, or (ii)~a demonstration that Marine-classified
``salient'' content is preferentially processed, remembered,
and valued by listeners in controlled conditions.

The corpus does not provide either (i)~or (ii)~with
the systematic controls that would establish the
correspondence. This is an instance of $\FailGeo{2}$
at moderate severity.

\section{Failure Geometry 3: Semantic Universality Collapse}
\label{sec:fg3-analysis}

\subsection{The Coherence Cascade}

The trajectory of the coherence concept traced in Chapter~9
provides the most systematic instance of $\FailGeo{3}$
in the Phoenix corpus. At each level of extension, the
concept's definition broadens to accommodate both positive
and negative instances: coherent signals are meaningful;
incoherent signals are degraded; coherence violations
are attention-capturing; coherence restoration is
therapeutic; consciousness is coherence; incoherence
is the pathology that consciousness pathologies share.

The accumulated effect of this cascade is a concept
that can accommodate any observation about salience,
attention, memory, and consciousness. This is the formal
structure of $\FailGeo{3}$: not that the concept is
wrong, but that it has been broadened until it no longer
generates specific predictions that could in principle
fail.

\begin{auditprop}
\label{ap:fg3-coherence}
The coherence concept in the Phoenix corpus exhibits
$\FailGeo{3}$ at high severity. The concept has been
extended across six distinct domains (signal processing,
auditory perception, memory, AI architecture, phenomenology,
consciousness) with successive broadening of its
operational definition at each step. At the broadest
extension, the concept can accommodate both the presence
and absence of any phenomenon in any domain by appeal
to ``coherence dynamics.'' The severity is high because
the collapse is domain-wide rather than local: it affects
all of the corpus's claims that involve the coherence
concept at the architectural and philosophical levels.
\end{auditprop}

\subsection{Scope and Salvageability}

The high severity of $\FailGeo{3}$ in the coherence
cascade does not make the entire Phoenix framework
unsalvageable. $\FailGeo{3}$ applies to the extended
uses of the coherence concept; it does not apply to
the engineering-level uses where the concept retains
its specific, operationalized definition. The Marine
jitter metric applied to audio signals is not exhibiting
$\FailGeo{3}$; it is exhibiting the tight, falsifiable
version of the coherence concept that the higher-level
extensions have lost.

The salvage strategy for $\FailGeo{3}$ is to identify
where the concept retains its specificity and to
restrict claims to those domains, while treating the
higher-level extensions as research questions rather
than established results. This is the strategy that
Part~IV will apply.

\section{Failure Geometry 4: Operational Severance}
\label{sec:fg4-analysis}

\subsection{The \texttt{.m8a} Format Revisited}

The consciousness metadata fields of the \texttt{.m8a}
format provide the clearest illustration of $\FailGeo{4}$.
The specific fields that exhibit Operational Severance are:

\textit{Emotional coordinates}: No measurement procedure
specified. No scale. No mapping from signal properties
or listener responses to coordinate values.

\textit{Resonance frequency}: Ambiguously defined. For
a physical resonator, the resonance frequency is
operationally accessible. For an audio file considered
as a whole, the operational definition is unclear.
Multiple non-equivalent operationalizations are possible.

\textit{Wave field coordinates}: These presuppose the
\MEMeight{} field, whose operational basis is itself
underspecified (see Chapter~7). A coordinate in an
unspecified field is doubly severed.

\textit{Audio DNA signature}: Named as the ``core
essence'' of the audio. No measurement procedure or
formal definition is provided.

\begin{auditprop}
\label{ap:fg4-m8a}
The \texttt{.m8a} format specification exhibits $\FailGeo{4}$
for the consciousness metadata fields at high severity.
These fields name quantities that would be meaningful
if measurable but for which no measurement procedure
has been specified. The severity is high because the
fields are central to the corpus's claim that the
\texttt{.m8a} format advances beyond conventional audio
formats in a specifically cognitive rather than merely
engineering sense. Without measurement procedures for
the consciousness metadata fields, this claim reduces
to a statement about the file's structure (it has fields
with these names) rather than about its content (these
fields contain meaningful measurements).
\end{auditprop}

\subsection{Severed Quantities in the MEM\textbar{}8 Formalism}

Beyond the \texttt{.m8a} format, the \MEMeight{} formalism
itself contains operationally severed quantities. The
most significant are the characteristic frequencies
$\omega_i$ and the representational space variable $x$.
As noted in Chapter~7, neither has been given a
measurement procedure that would allow their values
to be determined from observations.

The severity of this instance of $\FailGeo{4}$ is
moderate rather than high, because the \MEMeight{}
formalism's efficiency claims (the $O(N^2)$ implicit
associations at $O(N)$ storage cost) do not depend
on the operationalization of $x$ and $\omega_i$;
they follow from the mathematical structure of the
wave superposition, regardless of what the wave
functions represent. The unsalvageable component
of $\FailGeo{4}$ in \MEMeight{} is therefore the
consciousness claims that depend on the specific
physical or computational realization of the field,
not the associative memory claims that depend only
on its mathematical properties.

\section{Failure Geometry Interaction}
\label{sec:fg-interaction}

The four failure geometries do not operate independently.
They interact in ways that amplify each other's effects.

The most important interaction is between $\FailGeo{2}$
and $\FailGeo{3}$. When a concept is defined in terms
of the framework's primary variable ($\FailGeo{2}$),
the definition can be broadened without limit as the
primary variable is applied to new domains ($\FailGeo{3}$),
because the definition tracks the broadening automatically.
The authenticity concept exhibits this interaction:
because authenticity is implicitly defined as temporal
coherence, and because temporal coherence has been
extended to cover consciousness and all cognitive
phenomena, authenticity inherits the full scope of the
coherence cascade without requiring separate definitional
work.

The interaction between $\FailGeo{1}$ and $\FailGeo{4}$
is also significant. Decorative Formalism ($\FailGeo{1}$)
names quantities that appear specific but are
underspecified; Operational Severance ($\FailGeo{4}$)
names quantities that appear measurable but have no
measurement procedure. These can reinforce each other:
a formally well-organized framework expression
($\FailGeo{1}$ decoration) that uses operationally
severed quantities ($\FailGeo{4}$) produces claims
that appear doubly rigorous --- formally organized
and technically named --- but are doubly inaccessible.

\medskip

Chapter~13 turns from the failure geometry analysis
to the four-level taxonomy, examining how the Phoenix
corpus's claims distribute across the levels of
engineering utility, computational metaphor, cognitive
phenomenology, and ontological commitment. This
distribution determines which claims survive
decomposition in each of the four levels and which
require reinterpretation or withdrawal.
